\begin{lemma}[Existence of a minimizer of $J$ over the cone]
Let $F \in L^2(\gamma_n)$ be nonnegative $\gamma_n$-a.e., put
$m = \int F \, d\gamma_n$ and let $a > m$.  Then there exists
$u \in \mathcal{K}$ with
\[
  J(u) \leq J(v) \qquad\text{for every } v \in \mathcal{K},
\]
where $\mathcal{K}$ is the cone of \noderef{ObstacleCone} and $J$ the functional
of \noderef{ObstacleFunctional}.
\end{lemma}
\begin{proof}
By part (2) of \noderef{ObstacleCoercivity}, the infimum
$I = \inf_{v \in \mathcal{K}} J(v)$ is a finite real number with $I \leq 0$.

\emph{Step 1: a bounded minimizing sequence.}
By the definition of the infimum, for each integer $k \geq 1$ there is
$v_k \in \mathcal{K}$ with
\[
  I \leq J(v_k) < I + \tfrac1k .
\]
This uses only that $I$ is finite and does not presuppose that it is attained.
Since $I \leq 0$ we get $J(v_k) < I + 1 \leq 1$ for every $k$, so
\noderef{ObstacleSublevelBounded} applies with $C = 1$ and yields a finite
constant
$R^2 = M_1^2/(a-m)^2 + 8M_1$, $M_1 = 1 + \|F - m\mathbf{1}\|_2^2$, independent of
$k$, with
\begin{equation*}
  \|v_k\|_2^2 + \mathcal{E}(v_k) \leq R^2 \qquad\text{for every } k . \tag{1}
\end{equation*}
In particular $\|v_k\|_2 \leq R$ and $\mathcal{E}(v_k) \leq R^2$.  Also
$J(v_k) \to I$, since $I \leq J(v_k) < I + 1/k$.

\emph{Step 2: a weakly convergent subsequence in $L^2(\gamma_n)$.}
By (1) and the nonnegativity of $\mathcal{E}(v_k)$ (which holds by
\noderef{Energy}, the energy being a sum of nonnegative terms) we have
$\|v_k\|_2^2 \leq R^2$, hence $\|v_k\|_2 \leq R$ for every $k$, taking $R \geq 0$
to be the nonnegative square root of $R^2$.  Apply
\noderef{GaussL2WeakSeqCompact} to the sequence $(v_k)$ with this bound $R$: it
supplies indices $k_1 < k_2 < \cdots$ and $u \in L^2(\gamma_n)$ with
$\|u\|_2 \leq R$ and
\begin{equation*}
  \langle v_{k_j}, \varphi\rangle \longrightarrow \langle u, \varphi\rangle
  \qquad\text{for every } \varphi \in L^2(\gamma_n). \tag{2}
\end{equation*}
No compact embedding, and no compactness uniform in $n$, is used; in particular
no weak topology on $\mathcal{V}$ and no reflexivity of $\mathcal{V}$ is
invoked, which is where this proof departs from the reference paper's route.

\emph{Step 3: the limit is admissible.}
Each $v_{k_j}$ satisfies $v_{k_j} \geq 0$ $\gamma_n$-a.e.\ by
\noderef{ObstacleCone}, so \noderef{ConeWeaklyClosed}, applied to the sequence
$(v_{k_j})_j$ and the limit $u$ with hypothesis (2), gives $u \geq 0$
$\gamma_n$-a.e.  Moreover \noderef{EnergyWeaklyLSC}, applied to the same
sequence and limit, gives
\begin{equation*}
  \mathcal{E}(u) \leq \liminf_{j\to\infty} \mathcal{E}(v_{k_j}) \leq R^2 < \infty,
  \tag{3}
\end{equation*}
using (1).  By \noderef{ObstacleCone} these two facts say $u \in \mathcal{K}$.

\emph{Step 4: the linear term is weakly continuous.}
Since $\gamma_n$ is a probability measure (\noderef{GaussIsProb}), the constant
function $\mathbf{1}$ lies in $L^2(\gamma_n)$, hence so does
$a\mathbf{1} - F$.  Therefore, by (2) applied with
$\varphi = a\mathbf{1} - F$,
\begin{equation*}
  \int (a - F)\, v_{k_j}\, d\gamma_n
    = \langle v_{k_j},\, a\mathbf{1} - F\rangle
    \longrightarrow \langle u,\, a\mathbf{1} - F\rangle
    = \int (a-F)\, u \, d\gamma_n . \tag{4}
\end{equation*}

\emph{Step 5: $J(u) \leq I$.}
By \noderef{ObstacleFunctional} and then (3) and (4),
\[
  J(u) = \tfrac12 \mathcal{E}(u) + \int (a-F)u\, d\gamma_n
   \leq \tfrac12 \liminf_j \mathcal{E}(v_{k_j})
         + \lim_j \int (a-F)v_{k_j}\, d\gamma_n
   = \liminf_j J(v_{k_j}) ,
\]
where the last equality uses that the liminf of a sum of a sequence with a
liminf and a sequence with a limit is the sum of the liminf and the limit, all
quantities being finite reals by (1) and (3).  Since $J(v_k) \to I$ by Step 1, the
subsequence also satisfies $J(v_{k_j}) \to I$, so $\liminf_j J(v_{k_j}) = I$ and
hence $J(u) \leq I$.

\emph{Step 6: conclusion.}
By Step 3, $u \in \mathcal{K}$, so the definition of $I$ as the infimum over
$\mathcal{K}$ gives $I \leq J(u)$.  Combined with Step 5, $J(u) = I$, and
therefore $J(u) = I \leq J(v)$ for every $v \in \mathcal{K}$.
\end{proof}
